\section{Related Work}
\label{sec:related}

To systematically position our work within the current literature, we analyze 40 peer-reviewed studies organized across three core research themes.

\subsection{Theme 1: Traditional Machine Learning and Rule-Based Filtering}
Early research on telecommunication spam detection heavily relied on classical shallow algorithms, such as Support Vector Machines (SVM), Naive Bayes, and decision tree ensembles paired with TF-IDF or n-gram feature representations~\cite{tuan2023evaluating, xie2025comparing, ghourabi2023enhancing}. For instance, Saeed~\cite{saeed2023method} reported a $98.1\%$ accuracy using SVM on standard English benchmark datasets. However, Xie~\cite{xie2025comparing} demonstrated that shallow classifiers degrade significantly when encountering out-of-vocabulary teencode, deliberate orthographic obfuscations, and adversarial synonym replacements typical of modern Vietnamese cyber scams. Moreover, static rule-based filters fail to capture long-range semantic dependencies and subtle contextual urgency.

\subsection{Theme 2: Pretrained Language Models for Vietnamese NLP}
Transformer-based architectures have significantly improved natural language understanding in low-resource settings. In Vietnamese NLP, PhoBERT~\cite{nguyen2020phobert} established the primary benchmark for syllable-level contextual representations across downstream tasks. Building on pre-trained transformers, Nguyen-Xuan et al.~\cite{nguyenxuan2026mixed} introduced a mixed-language dual-encoder achieving an F1-score of $0.940$ on code-mixed phishing text, while Cam and Binh~\cite{cam2026vnsed} reached $86.47\%$ accuracy applying PhoBERT to Vietnamese email spam. Despite these advances, existing Vietnamese transformer models rely almost exclusively on standard symmetric cross-entropy loss. Because scam messages represent a minor fraction of real-world message traffic, unweighted loss functions bias decision thresholds toward the majority benign class, leading to high false-negative rates on subtle fraud lures.

\subsection{Theme 3: LLMs and Edge-Deployed Small Language Models}
Recent studies have evaluated Large Language Models (LLMs) and on-device Small Language Models (SLMs) for threat triage~\cite{saias2025advances, zhou2026fraudsmswalker, elzemity2026agentic}. Zhou et al.~\cite{zhou2026fraudsmswalker} benchmarked leading commercial models (GPT-4o, Claude, and Gemini) on SMS fraud, finding that scam recall varied from $64.16\%$ to $90.13\%$, while benign recall dropped as low as $12.81\%$ due to severe over-flagging. Conversely, lightweight mobile frameworks such as QuishingShield~\cite{banda2026quishingshield} and distilled edge models~\cite{mrinal2026resilient, phap2026edgeai} prioritize low-latency mobile inference, achieving execution times under $50\text{ ms}$ on CPUs. However, these edge architectures focus on English corpora and lack cost-sensitive penalty functions designed for the linguistic nuances of Vietnamese social media text.

\begin{table*}[t]
\centering
\caption{Comparative Overview of Representative Threat Detection Studies vs. Proposed ScamShield-VN}
\label{tab:related_work}
\begin{tabular}{lllllp{5.5cm}}
\toprule
\textbf{Study} & \textbf{Language} & \textbf{Model / Architecture} & \textbf{Evaluation Dataset} & \textbf{Key Metric} & \textbf{Primary Limitation / Research GAP} \\
\midrule
Tu\'{a}n et al.~\cite{tuan2023evaluating} & Vietnamese & PhoBERT, CNN, LSTM & Multi-accent SMS ($N\approx 3,000$) & Acc: $95.56\%$ & Evaluated symmetric loss; no cost-sensitive penalty. \\
Nguyen-Xuan~\cite{nguyenxuan2026mixed} & Viet-Eng & MLTEA Dual-Encoder & Code-mixed Phishing & F1: $0.940$ & Cloud-scale architecture; high latency for edge. \\
Cam \& Binh~\cite{cam2026vnsed} & Vietnamese & PhoBERT, BiLSTM & VNSED Email Corpus & Acc: $86.47\%$ & Lacks social slang vocabulary; high false negatives. \\
Zhou et al.~\cite{zhou2026fraudsmswalker} & Multilingual & GPT-4o, Claude, Gemini & FraudSMSWalker ($N\approx 4,000$) & Recall: $64.16\text{--}90.13\%$ & Extreme false positives (benign recall $<30.25\%$). \\
Banda et al.~\cite{banda2026quishingshield} & English & Llama-3.2-3B Mobile & QR/Text Lures ($N=4,200$) & Latency: $45\text{ ms}$ & English only; heavy RAM footprint ($>2\text{ GB}$). \\
Phap et al.~\cite{phap2026edgeai} & English & BiGRU + MaskedPool + WBCE & UCI \& LSDST ($N=5,574$) & Latency: $0.25\text{ ms}$ & Recurrent model; lacks transformer contextual semantics. \\
\midrule
\textbf{This Work} & \textbf{Vietnamese} & \textbf{ViSoBERT + WBCE ($\alpha=5.0$)} & \textbf{Merged SMS ($N=2,665$)} & \textbf{Recall: $95.62\%$} & \textbf{Pareto-optimal edge detection: $15\text{ ms}$, $390\text{ MB}$, INT8.} \\
\bottomrule
\end{tabular}
\end{table*}

\subsection{Positioning of ScamShield-VN}
Unlike existing approaches, ScamShield-VN directly addresses the trade-off between low-latency on-device execution and semantic accuracy for Vietnamese fraud detection. By combining a social-media-specialized transformer encoder (\texttt{uitnlp/visobert}) with weighted binary cross-entropy loss ($\mathcal{L}_{\text{WBCE}}$, $\alpha=5.0$) and a fail-safe multi-agent council, our architecture achieves high fraud recall with minimal on-device latency.

